
lemme facile suivant, donc pour \(X_{S'}/S'\) en vertu de 3.5 et 3.4 (ii), enfin pour \(X/S\) gr\^ace \`a 3.4 (i).

\textbf{Lemme 3.7.} \emph{Soient \(X=\coprod X_i\) de type fini sur \(S\) noeth\'erien et \(\mathcal F\) une famille de classes de faisceaux coh\'erents sur les fibres de \(X/S\). Pour que \(\mathcal F\) soit limit\'ee, il faut et il suffit que pour tout \(i\) la famille \(\mathcal F_i\) induite sur les fibres de \(X_i/S\) soit limit\'ee.}

Dans la d\'emonstration de 3.6, on s'est ramen\'e au cas o\`u \(X\) est projective et int\`egre sur \(S\)\,; munissons \(X/S\) d'un faisceau tr\`es ample \({O}_X(1)\). Soient \(\mathcal L\) une famille limit\'ee de classes de faisceaux inversibles sur les fibres de \(X/S\) et \(L\) un faisceau inversible sur un \(X_T\) avec \(T\) de type fini sur \(S\), qui limite \(\mathcal L\). Quitte \`a couper \(T\) en morceaux, on peut supposer \(L\) plat sur \(T\), et quitte \`a remplacer \(T\) par chacune de ses composantes connexes, on peut supposer \(T\) connexe. Alors le polyn\^ome en \(q\) et \(m\), \(P(q,m)=\chi(L_t^{\otimes q}(m))\) ne d\'epend pas de \(t\in T\) (EGA III 7.9.4).

Soit \(\mathcal L_n=\Phi_n^{-1}(\mathcal L)\)~: \(\mathcal L_n\) est la famille des classes des faisceaux inversibles \(M_K\) tels que la classe de \(M_K^{\otimes n}\) soit dans \(\mathcal L\). Soit \(\chi(M_K^{\otimes p}(m))=\sum_{i=0}^{r}a_i(p)\binom{m+i}{i}\), o\`u \(a_i(p)\) est un polyn\^ome en \(p\) de degr\'e \(\le r-i\). Pour \(q,m\) variables, on a \(\chi(M_K^{\otimes qn}(m))=P(q,m)\) et chaque \(a_i(qn)\) ne d\'epend donc pas du \(M_K\) choisi\,; par suite chaque polyn\^ome \(a_i(p)\) ne d\'epend pas non plus de \(M_K\). Par cons\'equent la famille \(\mathcal L_n\) a un seul polyn\^ome de Hilbert \(\chi(M_K(m))\), donc est limit\'ee en vertu de 2.11.

\textbf{Th\'eor\`eme 3.8.} \emph{Soient \(S\) un sch\'ema noeth\'erien, \(Y\) un \(S\)-sch\'ema projectif muni d'un faisceau inversible relativement ample \({O}_Y(1)\), \(X\) le sch\'ema des z\'eros d'une section \(\varphi\) de \({O}_Y(1)\), et \(f:X\to Y\) le morphisme d'inclusion. On}
